\section{Certification of the ensemble kinetic minimum}
\label{sec:ensemble}

We now pass from the weighted spectral estimate to the full fermionic
constrained search.  The argument uses the density regularity from
\Cref{lem:density-regularity}, the positive orbitals and their equations from
\Cref{prop:orbital-equations}, and \Cref{eq:weighted-gap}.  These hypotheses
will also suffice when they are established by the finite tests in
\Cref{sec:parameters}.

\subsection{Ensembles and the Pauli bound}

An ensemble is a positive trace-class operator $\Gamma$ on
$\mathfrak H_3=\bigwedge^3\mathfrak h$, where
$\mathfrak h=L^2(\R^3;\C^2)$, with $\Tr\Gamma=1$.  Its one-particle
reduced density matrix $\gamma_\Gamma$ is characterized by
\begin{equation}
 \Tr(\gamma_\Gamma A)
 =\Tr\left(\Gamma\sum_{i=1}^3 A_i\right)
 \qquad\text{for every bounded operator $A$ on $\mathfrak h$}.
 \label{eq:ensemble-rdm}
\end{equation}
In particular, $\gamma_\Gamma\ge0$ and $\Tr\gamma_\Gamma=3$.
For any unit spin orbital $f\in\mathfrak h$, the number operator
$\sum_i(|f\rangle\langle f|)_i$ is an orthogonal projection on
$\mathfrak H_3$.  Indeed, extend $f$ to an orthonormal basis of
$\mathfrak h$.  In the induced wedge basis, this operator has eigenvalue
one precisely on the determinants that contain $f$, and zero otherwise.
Its expectation is therefore at most one.  Applying
\Cref{eq:ensemble-rdm} gives
\begin{equation}
 0\le\gamma_\Gamma\le\id,
 \qquad \Tr\gamma_\Gamma=3.
 \label{eq:ensemble-pauli}
\end{equation}
This argument permits arbitrary complex spinors and arbitrary spin
coherences.

Write the spectral decomposition over positive eigenvalues as
$\gamma_\Gamma=\sum_j\nu_j|\psi_j\rangle\langle\psi_j|$.
The spin-summed density and kinetic expectation are
\begin{align}
 n_\Gamma(x)
 &=\sum_j\nu_j|\psi_j(x)|_{\C^2}^2,
 \label{eq:ensemble-density}\\
 T[\Gamma]
 &=\frac12\sum_j\nu_j\|\nabla\psi_j\|_2^2.
 \label{eq:ensemble-kinetic}
\end{align}
The second expression takes the value $+\infty$ if necessary; for finite
kinetic expectation every $\psi_j$ with $\nu_j>0$ belongs to
$H^1(\R^3;\C^2)$.  It equals the expectation of the three-electron
kinetic-energy form, by applying \Cref{eq:ensemble-rdm} to increasing
bounded spectral truncations of $-\Delta/2$ and using monotone
convergence.  The ensemble constrained-search value is
\begin{equation}
 T_{\mathrm{ens}}[n]
 =\inf\{T[\Gamma]:\Gamma\ge0,\ \Tr\Gamma=1,\ n_\Gamma=n\}.
 \label{eq:ensemble-search}
\end{equation}

Set $\widehat g=g/\sqrt2$, and let $P_g$ denote the orthogonal projection
onto the normalized spatial orbital $\widehat g$.  The orbitals
$\widehat g$ and $h$ are orthonormal by their norms and parity.  Thus
\begin{equation}
 \Phi_\theta
 = (\widehat g\uparrow)\wedge(\widehat g\downarrow)\wedge(h\uparrow)
 \label{eq:ensemble-determinant}
\end{equation}
is a normalized determinant, with finite kinetic energy and density
$2\widehat g^2+h^2=g^2+h^2=n$.

\begin{proposition}[An ensemble kinetic certificate]
\label{prop:ensemble-certificate}
Assume the density regularity of \Cref{lem:density-regularity}, the
orbital hypotheses of \Cref{prop:orbital-equations}, and the weighted
gap \Cref{eq:weighted-gap}.  Then the formal potential $V$ satisfies
$\int_{\R^3}n|V|<\infty$.  Every ensemble with density $n$ and finite
kinetic energy obeys
\begin{equation}
 T[\Gamma]+\int_{\R^3}Vn\ge d.
 \label{eq:ensemble-certificate}
\end{equation}
The determinant $\Phi_\theta$ attains equality.  Consequently,
\begin{equation}
 T_{\mathrm{ens}}[n]
 =T[|\Phi_\theta\rangle\langle\Phi_\theta|]
 =d-\int_{\R^3}Vn.
 \label{eq:ensemble-minimum}
\end{equation}
\end{proposition}

\begin{proof}
We first verify the absolute integrability required to use $V$.
By \Cref{lem:density-regularity}, $\Delta n\in L^1(\R^3)$ and
$\sqrt n\in H^1(\R^3)$.  Local positivity and regularity give the
almost-everywhere identity
\[
 \sqrt n\,\Delta\sqrt n
 =\frac12\Delta n-|\nabla\sqrt n|^2.
\]
Using the formula for $V$ in \Cref{prop:orbital-equations}, the finite
angle energy and $\int n\sin^2\theta=1$, we obtain
\begin{equation}
 \int n|V|
 \le \frac14\|\Delta n\|_1
      +\frac12\|\nabla\sqrt n\|_2^2
      +\frac12\int n|\nabla\theta|^2+d
 <\infty.
 \label{eq:ensemble-V-integrability}
\end{equation}
The local assertion $V\in L^2_{\mathrm{loc}}$ was proved in
\Cref{prop:orbital-equations}.

For a scalar orbital $f$, consider the quadratic expression
\[
 q_V[f]=\frac12\int|\nabla f|^2+\int V|f|^2,
 \qquad
 f\in\mathcal E_V
 :=\left\{f\in H^1(\R^3):\int|V||f|^2<\infty\right\}.
\]
Only this expression on its indicated test class is needed; we do not
assume an ordinary Hamiltonian or a closed form with potential $V$.
For $f\in C_c^\infty(\R^3;\C)$, the function $f/g$ is admissible in
\Cref{eq:weighted-gap}, since $g$ is positive and bounded below on each
compact set.  Testing $(-\Delta/2+V)g=0$ locally with $|f|^2/g$ gives
the positive-orbital identity
\begin{equation}
 q_V[f]
 =\frac12\int g^2\left|\nabla\left(\frac f g\right)\right|^2
 \ge d\left(\|f\|_2^2-|\langle\widehat g,f\rangle|^2\right).
 \label{eq:ensemble-orbital-bound}
\end{equation}
The identity uses real parts in the gradient expansion and therefore
holds for complex $f$ as well.
This is the positive ground-state representation; see
\cite[Section~2.3, equation~(2.12)]{FS08}.

We record the extension of \Cref{eq:ensemble-orbital-bound} carefully.
On a fixed compact set, $H^1$ approximation by smooth functions passes
through the potential integral: $V\in L^2$ there and strong $H^1$
convergence implies strong $L^4$ convergence, hence convergence of the
squared moduli in $L^2$.  Thus the inequality holds for compactly
supported $H^1$ orbitals.  Choose real cutoffs $0\le\chi_R\le1$ that
equal one on $|x|\le R$, vanish on $|x|\ge2R$, and satisfy
$|\nabla\chi_R|\le C/R$.  For $f\in\mathcal E_V$,
$\chi_Rf\to f$ in $H^1$, while dominated convergence gives
$\int V|\chi_Rf|^2\to\int V|f|^2$.  The projection term also converges
in $L^2$.  This proves \Cref{eq:ensemble-orbital-bound} on all of
$\mathcal E_V$.

Apply this scalar inequality to both components of a spinor $\psi$.
With $\widetilde P_g=P_g\otimes\id_{\C^2}$, the resulting estimate is
\begin{equation}
 \frac12\|\nabla\psi\|_2^2+\int V|\psi|_{\C^2}^2
 \ge d\bigl(\|\psi\|_2^2
             -\langle\psi,\widetilde P_g\psi\rangle\bigr).
 \label{eq:ensemble-spinor-bound}
\end{equation}
Now let $\Gamma$ have density $n$ and finite kinetic expectation.
By \Cref{eq:ensemble-density,eq:ensemble-V-integrability} and Tonelli's
theorem,
\[
 \sum_j\nu_j\int|V||\psi_j|_{\C^2}^2=\int|V|n<\infty.
\]
Hence every occupied natural spin orbital belongs componentwise to
$\mathcal E_V$, and the potential terms can be summed absolutely.
Summing \Cref{eq:ensemble-spinor-bound} and using
\Cref{eq:ensemble-pauli} gives
\begin{align*}
 T[\Gamma]+\int Vn
 &\ge d\left(\Tr\gamma_\Gamma
                -\Tr(\gamma_\Gamma\widetilde P_g)\right)\\
 &\ge d(3-2)=d.
\end{align*}
Here $\widetilde P_g$ has rank two and
$0\le\gamma_\Gamma\le\id$.  No rank bound on $\Gamma$ is used.

Finally, the orbitals $\widehat g,h$ belong to $\mathcal E_V$ by
\Cref{eq:ensemble-V-integrability}.  If $f$ denotes either of them,
with respective energy $\varepsilon=0,d$, testing its distributional
equation by $\chi_R^2 f$ gives
\[
 q_V[\chi_R f]
 =\varepsilon\|\chi_R f\|_2^2
   +\frac12\int|\nabla\chi_R|^2|f|^2.
\]
All terms converge as above, so $q_V[\widehat g]=0$ and $q_V[h]=d$.
The determinant in \Cref{eq:ensemble-determinant} therefore satisfies
\[
 T[|\Phi_\theta\rangle\langle\Phi_\theta|]+\int Vn
 =2q_V[\widehat g]+q_V[h]=d.
\]
This proves the minimum over all ensembles with finite kinetic energy;
ensembles with infinite kinetic expectation cannot lower it.
\end{proof}

\subsection{Identification of a representing potential}

\begin{proposition}[Transfer to a representing Hamiltonian]
\label{prop:potential-identification}
Assume the hypotheses of \Cref{prop:ensemble-certificate}.  Let
$v\in L^{3/2}(\R^3)+L^\infty(\R^3)$ be real and spin independent,
and let $h_v=-\Delta/2+v$ be its spatial one-particle operator, defined
by the closed form on $H^1(\R^3)$.  Suppose
$H_v=\sum_{i=1}^3(h_v\otimes\id_{\C^2})_i$ on $\mathfrak H_3$
has a ground eigenspace supporting an ensemble $\Gamma_s$ with
$n_{\Gamma_s}=n$.  Then $\Phi_\theta$ belongs to that ground eigenspace,
and there is a real constant $e_g$ such that
\begin{equation}
 h_v g=e_g g,
 \qquad v=V+e_g\quad\text{almost everywhere}.
 \label{eq:ensemble-potential-identification}
\end{equation}
In particular, $V\in L^{3/2}(\R^3)+L^\infty(\R^3)$.
\end{proposition}

\begin{proof}
The potential class supplies, for every $\varepsilon>0$, a constant
$C_\varepsilon$ such that
\begin{equation}
 \int|v||f|^2
 \le\varepsilon\,\frac12\|\nabla f\|_2^2
      +C_\varepsilon\|f\|_2^2,
 \qquad f\in H^1(\R^3).
 \label{eq:ensemble-standard-form-bound}
\end{equation}
Indeed, write $v=v_1+v_\infty$ and truncate $v_1$ at a sufficiently
large amplitude.  Its remaining $L^{3/2}$ tail has arbitrarily small
norm and is controlled by the Sobolev inequality; the truncated part
and $v_\infty$ are bounded.  Summing
\Cref{eq:ensemble-standard-form-bound} over spins and particles gives
\[
 q_{H_v}[\Xi]
 \ge\frac{1-\varepsilon}{2}\sum_{i=1}^3\|\nabla_i\Xi\|_2^2
       -3C_\varepsilon\|\Xi\|_2^2
\]
on the many-electron $H^1$ form domain.  If $E=\inf\sigma(H_v)$, every
normalized vector in its ground eigenspace consequently has kinetic
expectation at most $(E+3C_\varepsilon)/(1-\varepsilon)$ for any
$0<\varepsilon<1$.  Expanding $\Gamma_s$ in its trace-class spectral
decomposition proves that $T[\Gamma_s]<\infty$, even for an
infinite-dimensional ground eigenspace.  The absolute potential
expectation is finite by the same form bound, and
\[
 E=T[\Gamma_s]+\int vn.
\]

By \Cref{prop:ensemble-certificate}, the determinant has no larger
kinetic expectation and has the same density.  Thus
\[
 E\le q_{H_v}[\Phi_\theta]
 =T[|\Phi_\theta\rangle\langle\Phi_\theta|]+\int vn
 \le T[\Gamma_s]+\int vn=E.
\]
The square-root representation of the closed form
\cite[Chapter~VI, Theorem~2.23]{Kato} shows that equality at its bottom implies
$(H_v-E)^{1/2}\Phi_\theta=0$.  Hence $\Phi_\theta$ is itself a
ground-state eigenvector.

It remains to extract a spatial eigenfunction without any assumption
on the other ground states or on the multiplicities of one-particle
levels.  Denote the sesquilinear form of $h_v$ by
\[
 q_v(f_1,f_2)
 =\frac12\int\nabla\overline{f_1}\cdot\nabla f_2
       +\int v\overline{f_1}f_2,
\]
linear in its second argument.  For $f\in H^1(\R^3)$, replace the single down-spin orbital
in \Cref{eq:ensemble-determinant} and set
\[
 \Xi_f=(\widehat g\uparrow)\wedge(f\downarrow)\wedge(h\uparrow).
\]
This vector lies in the many-electron form domain.  Expanding the
one-body form between these determinants, spin orthogonality and
$\langle\widehat g,h\rangle=0$ give
\[
 \begin{aligned}
 q_{H_v}(\Phi_\theta,\Xi_f)
 &=q_v(\widehat g,f)\\
 &\quad+\langle\widehat g,f\rangle
       \bigl(q_v(\widehat g,\widehat g)+q_v(h,h)\bigr).
 \end{aligned}
\]
If $\langle\widehat g,f\rangle=0$, then $\Xi_f$ is orthogonal to
$\Phi_\theta$, and the ground-state form equation makes the left side
zero.  Therefore
\[
 q_v(\widehat g,f)=0
 \qquad\text{whenever }\langle\widehat g,f\rangle=0.
\]
Put $e_g=q_v(\widehat g,\widehat g)\in\R$.  Decomposing an arbitrary
$f\in H^1$ into its component parallel to $\widehat g$ and its
orthogonal component yields
\[
 q_v(\widehat g,f)=e_g\langle\widehat g,f\rangle
 \qquad(f\in H^1).
\]
By Hermitian symmetry this is equivalent to
\(q_v(f,\widehat g)=e_g\langle f,\widehat g\rangle\).
The closed-form representation theorem
\cite[Chapter~VI, Theorem~2.1\textup{(iii)}]{Kato} therefore gives
$\widehat g\in\Dom(h_v)$ and $h_v\widehat g=e_g\widehat g$.

Subtracting $(-\Delta/2+V)g=0$ now gives
$(v-V-e_g)g=0$ in distributions.  The product is locally integrable:
$v\in L^{3/2}_{\mathrm{loc}}$, $V\in L^2_{\mathrm{loc}}$, and
$g\in H^2_{\mathrm{loc}}$ is locally bounded.  Since $g$ is continuous
and strictly positive, division on each compact set proves
$v=V+e_g$ almost everywhere.  Adding a constant preserves the stated
potential class, completing the proof.
\end{proof}
